\section{Methodology}

We present \textbf{Flux \& Key}, a training-free, reference-guided image editing pipeline
for multi-modal diffusion transformers.
Given a source image $\mathbf{I}_A \in \mathbb{R}^{H \times W \times 3}$, a reference image
$\mathbf{I}_B \in \mathbb{R}^{H' \times W' \times 3}$ containing the specific object to be
transferred, a binary spatial mask $\mathbf{M}_A \in \{0,1\}^{H \times W}$ delineating the
edit region, and a text prompt $\mathcal{T}$, our goal is to produce an edited image
$\hat{\mathbf{I}}$ satisfying three constraints simultaneously:
(i)~the visual identity of the specific object in $\mathbf{I}_B$ is faithfully transferred
into the masked region;
(ii)~all pixels outside $\mathbf{M}_A$ are reconstructed with high fidelity; and
(iii)~no model fine-tuning or iterative ODE inversion is required.
The pipeline is summarised in Figure~\ref{fig:pipeline} and detailed below.

% ---------------------------------------------------------------------------
\subsection{Reference Object Extraction}
% ---------------------------------------------------------------------------

We first isolate the target object in $\mathbf{I}_B$ using Grounding
DINO~\cite{liu2023grounding}, which localises the object via text-conditioned
bounding-box detection.
The detected bounding box defines the crop region:
\begin{equation}
    \mathbf{I}_B^{\text{crop}} = \operatorname{Crop}\!\left(\mathbf{I}_B,\;
        \operatorname{bbox}_{\text{DINO}}(\mathbf{I}_B, \mathcal{T})\right).
    \label{eq:crop}
\end{equation}
To address the scale mismatch between the reference object and the target region, we resize
$\mathbf{I}_B^{\text{crop}}$ to match the bounding-box extent of $\mathbf{M}_A$:
\begin{equation}
    \mathbf{I}_B^{\text{aligned}} =
        \operatorname{Resize}\!\left(
            \mathbf{I}_B^{\text{crop}},\;
            \operatorname{bbox}(\mathbf{M}_A)
        \right).
    \label{eq:resize}
\end{equation}
We then encode $\mathbf{I}_B^{\text{aligned}}$ through the VAE encoder $\mathcal{E}$,
preserving spatial and spectral structure:
\begin{equation}
    \mathbf{z}_B = \mathcal{E}\!\left(\mathbf{I}_B^{\text{aligned}}\right),
    \label{eq:ref_vae}
\end{equation}
and similarly encode the source:
\begin{equation}
    \mathbf{z}_A = \mathcal{E}(\mathbf{I}_A).
    \label{eq:src_vae}
\end{equation}
Note that we deliberately avoid encoding $\mathbf{I}_B$ through CLIP, which is optimised
for semantic similarity and discards the high-frequency instance detail — exact texture,
surface finish, edge geometry — that distinguishes a specific object from a generic
category~\cite{radford2021clip}.

% ---------------------------------------------------------------------------
\subsection{Dual-Stream Feature Extraction (Two Clean Forward Passes)}
% ---------------------------------------------------------------------------

In order to obtain transformer-level attention features from both $\mathbf{I}_A$ and
$\mathbf{I}_B^{\text{aligned}}$, we perform two \emph{clean} (noise-free) forward passes
through the frozen FLUX transformer $f_\theta$ at $t = 0$.
No inversion is performed; these are standard feed-forward evaluations on clean latents.

\paragraph{Source stream.}
A forward pass on $\mathbf{z}_A$ yields background key-value features for all $L$ blocks:
\begin{equation}
    \left\{
        \mathbf{K}_{\mathrm{src}}^{(\ell)},\;
        \mathbf{V}_{\mathrm{src}}^{(\ell)}
    \right\}_{\ell=1}^{L}
    = f_\theta(\mathbf{z}_A)\big|_{t=0}.
    \label{eq:src_feat}
\end{equation}

\paragraph{Reference stream.}
A forward pass on $\mathbf{z}_B$ yields reference key-value features:
\begin{equation}
    \left\{
        \mathbf{K}_{\mathrm{ref}}^{(\ell)},\;
        \mathbf{V}_{\mathrm{ref}}^{(\ell)}
    \right\}_{\ell=1}^{L}
    = f_\theta(\mathbf{z}_B)\big|_{t=0}.
    \label{eq:ref_feat}
\end{equation}
In principle, the 2D rotary position embeddings (RoPE) of the reference tokens could be
re-indexed to the spatial coordinates of the foreground region
$\mathcal{F} = \{i : \mathbf{M}_A^{(i)} = 1\}$ within Image~A's latent grid, so that
reference features carry position context aligned to the insertion site.
In our current implementation we use standard RoPE coordinates; RoPE re-indexing is left
as a planned improvement (see Ablation Plan, Section~\ref{sec:ablation}).
The total cost of Stages~1--2 is two forward passes (equivalent to two inversion steps),
after which no further passes on either source or reference are needed.

% ---------------------------------------------------------------------------
\subsection{Selective Stochastic Inversion (SSI)}
% ---------------------------------------------------------------------------

We need to initialise the denoising trajectory differently for foreground and background
tokens.
Background tokens should start from their clean latent values so that the frequency lock
(Section~\ref{sec:bg_lock}) can exactly reproduce the source.
Foreground tokens need sufficient freedom to regenerate a structurally different object.

\textbf{FLUX uses rectified flow}, not the DDPM noise schedule.
The forward noising process is:
\begin{equation}
    \mathbf{z}_t = (1 - t)\,\mathbf{z}_0 + t\,\boldsymbol{\varepsilon},
    \quad \boldsymbol{\varepsilon} \sim \mathcal{N}(\mathbf{0}, \mathbf{I}),
    \quad t \in [0,1].
    \label{eq:flow_noise}
\end{equation}
We apply SSI exclusively to the foreground token subset
$\mathbf{z}_A^{\mathcal{F}}$ at a chosen stochasticity level $\tau$:
\begin{align}
    \tilde{\mathbf{z}}_\tau^{\mathcal{F}}
        &= (1-\tau)\,\mathbf{z}_A^{\mathcal{F}} + \tau\,\boldsymbol{\varepsilon},
    \label{eq:ssi} \\
    \tilde{\mathbf{z}}_0^{\mathcal{B}}
        &= \mathbf{z}_A^{\mathcal{B}},
    \label{eq:bg_clean}
\end{align}
where $\tau \in [0,1]$ controls edit aggressiveness.
Background tokens are kept clean ($\tau=0$) and are never perturbed.
Unlike FSI-Edit~\cite{yang2026fsi}, which injects noise globally, SSI is spatially confined
to $\mathcal{F}$, eliminating mask-leakage artefacts by construction.
Denoising then begins from timestep $\tau$ for the foreground and from $t=0$
for the background (which the frequency lock will handle).

The stochasticity level $\tau$ is calibrated per edit category:
\begin{itemize}[nosep]
    \item Wholesale replacement (change object, add object): $\tau \in [0.7, 0.9]$
    \item Attribute modification (colour, texture, pose): $\tau \in [0.3, 0.5]$
\end{itemize}

% ---------------------------------------------------------------------------
\subsection{Inversion-Free Background Preservation via Frequency Locking}
\label{sec:bg_lock}
% ---------------------------------------------------------------------------

For background tokens $\mathcal{B} = \{i : \mathbf{M}_A^{(i)} = 0\}$, we apply
FIA-Edit-style frequency fusion~\cite{yang2026fia} \emph{masked to $\mathcal{B}$},
replacing KV-Edit's hard inversion-based cache with a soft frequency lock.

We define complementary high- and low-pass spectral filters.
For a $512\times512$ image, the FLUX VAE produces $64\times64$ latents (8$\times$
downsampling); $2\times2$ patch packing reduces this to a $32\times32 = 1024$-token sequence.
We apply the frequency lock along this token dimension.
Let $\mathcal{F}_{1D}$ denote the 1-D DFT along the token sequence axis and $\gamma$ the
fraction of low-frequency bins assigned to the target:
\begin{equation}
    \mathcal{H}(\hat{\mathbf{X}})_{k}
        = \hat{\mathbf{X}}_{k} \cdot \mathbf{1}[k > \lfloor\gamma K\rfloor],
    \quad
    \mathcal{L}(\hat{\mathbf{X}})_{k}
        = \hat{\mathbf{X}}_{k} \cdot \mathbf{1}[k \leq \lfloor\gamma K\rfloor],
    \label{eq:filters}
\end{equation}
where $K$ is the total number of frequency bins.
At transformer block $\ell$, during denoising at timestep $t$, background keys and values
are constructed as:
\begin{align}
    \mathbf{K}_{\mathcal{B}}^{(\ell)}
        &= \mathcal{F}_{1D}^{-1}\!\!\left(
              \mathcal{H}\!\left(\mathcal{F}_{1D}(\mathbf{K}_{\mathrm{src}}^{(\ell)})\right)
            + \mathcal{L}\!\left(\mathcal{F}_{1D}(\mathbf{K}_{\mathrm{tgt}}^{(\ell)})\right)
           \right)\bigg|_{\mathcal{B}},
    \label{eq:bg_K} \\
    \mathbf{V}_{\mathcal{B}}^{(\ell)}
        &= \mathcal{F}_{1D}^{-1}\!\!\left(
              \mathcal{H}\!\left(\mathcal{F}_{1D}(\mathbf{V}_{\mathrm{src}}^{(\ell)})\right)
            + \mathcal{L}\!\left(\mathcal{F}_{1D}(\mathbf{V}_{\mathrm{tgt}}^{(\ell)})\right)
           \right)\bigg|_{\mathcal{B}},
    \label{eq:bg_V}
\end{align}
where $\mathbf{K}_{\mathrm{tgt}}^{(\ell)}, \mathbf{V}_{\mathrm{tgt}}^{(\ell)}$ are computed
from the current noisy latent $\mathbf{z}_t$ at each denoising step.
Before mixing, we normalise $\mathbf{K}_{\mathrm{src}}$ to match the per-token
$\ell_2$ norm of $\mathbf{K}_{\mathrm{tgt}}$, preventing phase-amplitude interference
caused by the different noise levels of the two signals.

We also apply a time-aware schedule to $\gamma$: at high noise levels the source and target
features are at very different scales, so we let $\gamma$ approach $1$ (target-only) and
tighten the lock gradually as denoising progresses:
\begin{equation}
    \gamma_t = \gamma + (1 - \gamma)\,t,
    \label{eq:gamma_schedule}
\end{equation}
where $t \in [0, 1]$ is the current denoising timestep.
At $t \approx 1$, $\gamma_t \approx 1$ (no source injection); at $t = 0$, $\gamma_t = \gamma$
(full frequency lock).

\textit{Rationale.}
High-frequency source components encode edge geometry, texture, and fine spatial structure
of the background — the signals that define its visual identity.
Low-frequency target components carry global illumination and semantic context, permitting
subtle lighting adaptation consistent with the inserted object.
Confining the fusion to $\mathcal{B}$ (rather than the full image as in FIA-Edit) prevents
reference-object features from contaminating background tokens.

% ---------------------------------------------------------------------------
\subsection{Reference-Conditioned Foreground Injection}
% ---------------------------------------------------------------------------

For foreground tokens $\mathcal{F}$, we replace the text-only KV stream of KV-Edit
with a soft interpolation between the text branch and the reference branch:
\begin{align}
    \tilde{\mathbf{K}}_{\mathcal{F}}^{(\ell)}
        &= \alpha^{(\ell)}\,\mathbf{K}_{\mathrm{text}}^{(\ell)}
         + \left(1 - \alpha^{(\ell)}\right)\mathbf{K}_{\mathrm{ref}}^{(\ell)},
    \label{eq:fg_K} \\
    \tilde{\mathbf{V}}_{\mathcal{F}}^{(\ell)}
        &= \beta^{(\ell)}\,\mathbf{V}_{\mathrm{text}}^{(\ell)}
         + \left(1 - \beta^{(\ell)}\right)\mathbf{V}_{\mathrm{ref}}^{(\ell)},
    \label{eq:fg_V}
\end{align}
where $\mathbf{K}_{\mathrm{text}}^{(\ell)}, \mathbf{V}_{\mathrm{text}}^{(\ell)}$ are
produced from the text embedding $\mathcal{T}$, and
$\alpha^{(\ell)}, \beta^{(\ell)} \in [0,1]$ are layer-wise mixing coefficients.

\paragraph{Layer-wise schedule.}
Early transformer blocks govern global layout and pose; late blocks refine surface
appearance and fine detail~\cite{yang2026fia}.
We adopt a linearly decaying schedule that makes early blocks text-dominant
(preserving spatial structure from $\mathcal{T}$) and late blocks reference-dominant
(sharpening instance-specific detail from $\mathbf{I}_B$):
\begin{equation}
    \alpha^{(\ell)}
        = \alpha_{\max} - \frac{\ell}{L}\left(\alpha_{\max} - \alpha_{\min}\right),
    \label{eq:schedule}
\end{equation}
with $\alpha_{\max} = 0.8$, $\alpha_{\min} = 0.2$; $\beta^{(\ell)}$ follows the same
schedule.
These values are initial estimates subject to ablation (Section~\ref{sec:ablation}).

The per-block attention for foreground tokens is:
\begin{equation}
    \mathbf{A}_{\mathcal{F}}^{(\ell)}
        = \operatorname{Softmax}\!\!\left(
              \frac{\mathbf{Q}_{\mathcal{F}}^{(\ell)}\,
                    \tilde{\mathbf{K}}_{\mathcal{F}}^{(\ell)\top}}{\sqrt{d}}
          \right)\tilde{\mathbf{V}}_{\mathcal{F}}^{(\ell)}.
    \label{eq:fg_attn}
\end{equation}
Background and foreground attention streams are disjoint: background tokens never attend
to foreground or reference features, and vice versa.

\paragraph{Remark on clean vs.\ noisy reference features.}
Both $\mathbf{K}_{\mathrm{src}}^{(\ell)}$ and $\mathbf{K}_{\mathrm{ref}}^{(\ell)}$ are
computed from clean latents ($t=0$), while the query and target features during denoising
are computed from noisy latents at timestep $t$.
This domain mismatch is shared with FIA-Edit, which argues that high-frequency structural
signals are largely invariant to the noise level.
For the foreground interpolation we mitigate this by conditioning the mixing weight
$\alpha^{(\ell)}$ on $t$: at early timesteps (large $t$) the reference weight is
attenuated to avoid imposing clean-domain structure prematurely, with full reference
injection at $t \approx 0$:
\begin{equation}
    \alpha^{(\ell)}_t = \alpha^{(\ell)} + (1 - \alpha^{(\ell)})\,t.
    \label{eq:time_schedule}
\end{equation}

% ---------------------------------------------------------------------------
\subsection{Inference Algorithm}
% ---------------------------------------------------------------------------

The complete pipeline requires no ODE inversion and proceeds as follows:

\begin{enumerate}[leftmargin=*, nosep]
    \item \textbf{Extract reference:} Grounded-DINO + SAM~2 $\to$
          $\mathbf{I}_B^{\text{aligned}}$.
    \item \textbf{Encode:} $\mathbf{z}_A = \mathcal{E}(\mathbf{I}_A)$;\quad
          $\mathbf{z}_B = \mathcal{E}(\mathbf{I}_B^{\text{aligned}})$.
    \item \textbf{Source forward pass} (clean, $t=0$):
          extract $\{\mathbf{K}_{\mathrm{src}}^{(\ell)}, \mathbf{V}_{\mathrm{src}}^{(\ell)}\}$.
    \item \textbf{Reference forward pass} (clean, $t=0$):
          extract $\{\mathbf{K}_{\mathrm{ref}}^{(\ell)}, \mathbf{V}_{\mathrm{ref}}^{(\ell)}\}$.
    \item \textbf{SSI:}
          $\tilde{\mathbf{z}}_\tau^{\mathcal{F}} = (1-\tau)\mathbf{z}_A^{\mathcal{F}}
          + \tau\boldsymbol{\varepsilon}$;\quad
          $\tilde{\mathbf{z}}_0^{\mathcal{B}} = \mathbf{z}_A^{\mathcal{B}}$.
    \item \textbf{Denoising} from $t=\tau \to 0$: at each step and block $\ell$,
          background tokens use frequency-locked KVs
          (Eqs.~\ref{eq:bg_K}--\ref{eq:bg_V});
          foreground tokens use reference-interpolated KVs
          (Eqs.~\ref{eq:fg_K}--\ref{eq:fg_attn}).
    \item \textbf{Decode:} $\hat{\mathbf{I}} = \mathcal{D}(\mathbf{z}_{\text{final}})$.
\end{enumerate}

\paragraph{Computational complexity.}
The method adds two clean forward passes (Steps~3--4) relative to standard text-guided
inference, plus $\mathcal{O}(N \log N)$ FFT operations per block per denoising step
(where $N = 32 \times 32 = 1024$ tokens after $2\times2$ patch packing).
No ODE inversion is performed at any stage.

% ---------------------------------------------------------------------------
\subsection{Ablation Plan}
\label{sec:ablation}
% ---------------------------------------------------------------------------

To isolate the contribution of each component we plan the following ablations:

\begin{enumerate}[nosep]
    \item \textbf{Background lock:} hard KV-cache (KV-Edit) vs.\ frequency lock (ours).
    \item \textbf{Foreground conditioning:} text-only vs.\ reference-only
          ($\alpha=0$) vs.\ interpolation (ours, $\alpha=0.5$) vs.\ full schedule.
    \item \textbf{SSI level:} $\tau \in \{0.3, 0.5, 0.7, 0.9\}$.
    \item \textbf{Mixing schedule:} constant $\alpha$ vs.\ layer-wise
          (Eq.~\ref{eq:schedule}) vs.\ layer+time-aware
          (Eq.~\ref{eq:time_schedule}).
    \item \textbf{Reference encoding:} CLIP~\cite{radford2021clip}
          vs.\ DINOv2~\cite{oquab2023dinov2} vs.\ FLUX-VAE + transformer (ours).
    \item \textbf{RoPE re-indexing:} standard coordinates vs.\ re-indexed to
          foreground mask positions.
    \item \textbf{Frequency filter dimensionality:} 1-D token-sequence DFT (current)
          vs.\ 2-D spatial DFT on the $32\times32$ latent grid.
\end{enumerate}
